        公开研究情绪集中在三点。第一，WF6供应紧张叙事把中船特气、中巨芯、昊华科技等放到同一张表中；第二，钨价上涨强化了上游钨粉成本和战略资源安全的想象；第三，氟化工公司被同时赋予制冷剂配额、电子级HF、氟化液和含氟电子特气多重映射。

        但公开澄清也很明确：中巨芯提示高纯WF6尚未签署新的长期或大额实质性订单，昊华科技披露WF6收入占比低，华特气体提示无WF6合成产能且纯化业务占比较小，中船特气虽有产能和订单洽谈转述但未公开披露合同金额和WF6价格。我们的估值因此不把市场传闻价格直接当作业绩，而是单独设置市场情绪锚。

        \begin{riskbox}[共识拥挤风险]
        如果后续公告没有给出订单、价格、客户认证或扩产兑现，电子特气龙头的当前估值更容易从“国产替代稀缺资产”切换为“单品传闻驱动的高位拥挤交易”。
        \end{riskbox}


\begin{exhibitbox}[市场盘口与日内风险（1）]
\scriptsize
\begin{tabularx}{\exhibitboxwidth}{L{0.9cm} L{1.15cm} L{0.9cm} L{0.9cm} L{0.9cm} L{1.2cm} X X}
\toprule
\textbf{代码} & \textbf{名称} & \textbf{现价} & \textbf{涨跌幅} & \textbf{日内振幅} & \textbf{成交额} & \textbf{盘口状态} & \textbf{估值用途} \\
\midrule
002842 & 翔鹭钨业 & 47.53 & -8.9\% & 11.4\% & 19.0亿元 & 单日大跌；日内位置缺失；中高成交额 & 风险释放/验证公告支撑，降低追价权重 \\
603379 & 三美股份 & 66.83 & 1.1\% & 10.2\% & 11.0亿元 & 温和波动；日内位置缺失；中高成交额 & 流动性和价格信号作为监测变量 \\
300037 & 新宙邦 & 82.43 & -1.2\% & 7.5\% & 17.8亿元 & 温和波动；日内位置缺失；中高成交额；有结构化事件 & 流动性和价格信号作为监测变量 \\
605020 & 永和股份 & 39.91 & -3.2\% & 11.4\% & 15.8亿元 & 下跌；日内位置缺失；中高成交额；有结构化事件 & 流动性和价格信号作为监测变量 \\
002378 & 章源钨业 & 33.78 & -6.7\% & 8.7\% & 17.4亿元 & 单日大跌；日内位置缺失；中高成交额；有结构化事件 & 风险释放/验证公告支撑，降低追价权重 \\
600160 & 巨化股份 & 50.72 & -0.5\% & 13.1\% & 51.0亿元 & 温和波动；日内位置缺失；高成交额 & 主题交易活跃，需用订单/价格/毛利维持市场锚 \\
600549 & 厦门钨业 & 81.03 & -4.1\% & 8.2\% & 53.9亿元 & 下跌；日内位置缺失；高成交额 & 主题交易活跃，需用订单/价格/毛利维持市场锚 \\
002407 & 多氟多 & 45.66 & 10.0\% & 0.0\% & 15.9亿元 & 单日大涨；日内位置缺失；中高成交额；有结构化事件 & 流动性和价格信号作为监测变量 \\
600378 & 昊华科技 & 60.13 & -3.5\% & 12.7\% & 34.1亿元 & 下跌；日内位置缺失；中高成交额；有结构化事件 & 流动性和价格信号作为监测变量 \\
\bottomrule
\end{tabularx}
\end{exhibitbox}

\begin{exhibitbox}[市场盘口与日内风险（2）]
\scriptsize
\begin{tabularx}{\exhibitboxwidth}{L{0.9cm} L{1.15cm} L{0.9cm} L{0.9cm} L{0.9cm} L{1.2cm} X X}
\toprule
\textbf{代码} & \textbf{名称} & \textbf{现价} & \textbf{涨跌幅} & \textbf{日内振幅} & \textbf{成交额} & \textbf{盘口状态} & \textbf{估值用途} \\
\midrule
300346 & 南大光电 & 86.27 & 10.4\% & 15.8\% & 90.7亿元 & 单日大涨；日内位置缺失；高成交额；有结构化事件 & 拥挤升温，只提高市场锚观察，不提高基本面锚 \\
000657 & 中钨高新 & 93.04 & -6.5\% & 9.9\% & 67.5亿元 & 单日大跌；日内位置缺失；高成交额；有结构化事件 & 风险释放/验证公告支撑，降低追价权重 \\
688146 & 中船特气 & 339.00 & -12.1\% & 21.0\% & 54.0亿元 & 单日大跌；日内位置缺失；高成交额 & 风险释放/验证公告支撑，降低追价权重 \\
688549 & 中巨芯-U & 34.43 & 7.3\% & 10.8\% & 31.8亿元 & 单日大涨；日内位置缺失；中高成交额；有结构化事件 & 流动性和价格信号作为监测变量 \\
603505 & 金石资源 & 18.85 & -0.8\% & 10.9\% & 4.8亿元 & 温和波动；日内位置缺失；成交一般；有结构化事件 & 流动性和价格信号作为监测变量 \\
688268 & 华特气体 & 240.90 & 3.8\% & 10.7\% & 19.1亿元 & 上涨；日内位置缺失；中高成交额；有结构化事件 & 流动性和价格信号作为监测变量 \\
002971 & 和远气体 & 60.93 & -7.8\% & 16.6\% & 14.3亿元 & 单日大跌；日内位置缺失；中高成交额；有结构化事件 & 风险释放/验证公告支撑，降低追价权重 \\
002549 & 凯美特气 & 20.89 & 10.0\% & 3.0\% & 12.2亿元 & 单日大涨；日内位置缺失；中高成交额；有结构化事件 & 流动性和价格信号作为监测变量 \\
688106 & 金宏气体 & 37.76 & 9.4\% & 11.2\% & 20.2亿元 & 单日大涨；日内位置缺失；中高成交额 & 流动性和价格信号作为监测变量 \\
\bottomrule
\end{tabularx}
\sourcenote{本地实时行情能力层派生；盘口、涨跌幅、日内振幅和成交额只用于二级市场拥挤度和市场锚监测，不替代客户订单、合同平均售价、收入确认比例或单品毛利。}
\end{exhibitbox}

盘口表补上了“当前价格之外发生了什么”，但此前仍缺一个关键二级市场口径：官方现价、官方总市值、流通市值、换手率和量比。下面的东方财富官方行情快照已经进入本报告估值现价、市值、市盈率/市销率/市净率和空间测算，本地行情能力层价格仅保留为对照；但换手率、量比和成交额仍只用于自由流通盘约束与交易拥挤度审计，不能替代订单、合同平均售价或单品毛利证据。


\begin{exhibitbox}[官方市值与流动性口径（1）]
\scriptsize
\begin{tabularx}{\exhibitboxwidth}{L{0.85cm} L{1.15cm} L{1.2cm} L{1.2cm} L{0.9cm} L{0.9cm} L{0.8cm} X}
\toprule
\textbf{代码} & \textbf{名称} & \textbf{总市值} & \textbf{流通市值} & \textbf{流通/总} & \textbf{换手率} & \textbf{量比} & \textbf{流动性信号} \\
\midrule
600549 & 厦门钨业 & 1286.4亿元 & 1259.2亿元 & 97.9\% & 4.1\% & 1.19 & 高成交额；中高换手；流通盘约束中性 \\
000657 & 中钨高新 & 2120.0亿元 & 1351.8亿元 & 63.8\% & 4.8\% & 1.31 & 高成交额；中高换手；流通盘约束中性 \\
002378 & 章源钨业 & 405.8亿元 & 403.7亿元 & 99.5\% & 4.2\% & 0.94 & 中高成交额；中高换手；流通盘约束中性 \\
002842 & 翔鹭钨业 & 155.5亿元 & 127.6亿元 & 82.1\% & 14.2\% & 1.17 & 中高成交额；高换手；流通盘约束中性 \\
603505 & 金石资源 & 158.7亿元 & 158.7亿元 & 100.0\% & 2.9\% & 1.59 & 成交一般；换手温和；流通盘约束中性 \\
600160 & 巨化股份 & 1369.3亿元 & 1369.3亿元 & 100.0\% & 3.7\% & 1.67 & 高成交额；换手温和；流通盘约束中性 \\
603379 & 三美股份 & 408.0亿元 & 408.0亿元 & 100.0\% & 2.6\% & 1.98 & 中高成交额；换手温和；流通盘约束中性 \\
605020 & 永和股份 & 203.9亿元 & 200.7亿元 & 98.5\% & 7.6\% & 1.54 & 中高成交额；中高换手；流通盘约束中性 \\
002407 & 多氟多 & 543.6亿元 & 492.6亿元 & 90.6\% & 3.2\% & 0.33 & 中高成交额；换手温和；流通盘约束中性 \\
\bottomrule
\end{tabularx}
\end{exhibitbox}

\begin{exhibitbox}[官方市值与流动性口径（2）]
\scriptsize
\begin{tabularx}{\exhibitboxwidth}{L{0.85cm} L{1.15cm} L{1.2cm} L{1.2cm} L{0.9cm} L{0.9cm} L{0.8cm} X}
\toprule
\textbf{代码} & \textbf{名称} & \textbf{总市值} & \textbf{流通市值} & \textbf{流通/总} & \textbf{换手率} & \textbf{量比} & \textbf{流动性信号} \\
\midrule
600378 & 昊华科技 & 775.7亿元 & 644.9亿元 & 83.1\% & 5.0\% & 1.28 & 中高成交额；中高换手；流通盘约束中性 \\
300037 & 新宙邦 & 621.4亿元 & 445.8亿元 & 71.7\% & 3.9\% & 1.19 & 中高成交额；换手温和；流通盘约束中性 \\
688146 & 中船特气 & 1794.7亿元 & 491.5亿元 & 27.4\% & 10.6\% & 1.64 & 高成交额；高换手；流通盘约束强 \\
688549 & 中巨芯-U & 508.6亿元 & 202.9亿元 & 39.9\% & 15.7\% & 1.84 & 中高成交额；高换手；流通盘约束强 \\
688268 & 华特气体 & 307.6亿元 & 307.6亿元 & 100.0\% & 6.1\% & 1.07 & 中高成交额；中高换手；流通盘约束中性 \\
300346 & 南大光电 & 596.3亿元 & 565.9亿元 & 94.9\% & 15.8\% & 1.92 & 高成交额；高换手；流通盘约束中性 \\
688106 & 金宏气体 & 202.2亿元 & 202.2亿元 & 100.0\% & 9.9\% & 1.91 & 中高成交额；高换手；流通盘约束中性 \\
002549 & 凯美特气 & 145.3亿元 & 144.6亿元 & 99.6\% & 8.5\% & 1.47 & 中高成交额；高换手；流通盘约束中性 \\
002971 & 和远气体 & 128.3亿元 & 95.3亿元 & 74.3\% & 13.8\% & 1.60 & 中高成交额；高换手；流通盘约束中性 \\
\bottomrule
\end{tabularx}
\sourcenote{东方财富push2批量行情接口；官方现价和官方总市值进入估值现价、市值、市盈率/市销率/市净率和空间测算，流通市值、换手率和量比用于二级市场流通盘约束和拥挤交易审计；不替代客户订单、合同平均售价、收入确认比例或单品毛利。完整快照见 data/official\_market\_liquidity\_snapshot.json/md。}
\end{exhibitbox}

这组表把“钱在交易什么”拆得更细：总市值约束估值排序，流通市值和换手率约束拥挤度，量比约束当日成交放大程度。它仍然不能证明WF6订单、合同价格或单品毛利已经兑现；真正能改变基本面目标价的，仍是订单金额、公司合同平均售价、收入确认比例和单品毛利。


有了官方市值和换手率，还要看谁持有这些流通筹码。下面的股东结构与机构持仓快照把第一大股东、前十股东、流通前十和基金持仓放到同一张二级市场门禁表中，用来判断筹码集中度、机构参与度和追价容忍度。它只能解释“价格为什么容易被推高或快速回撤”，不能证明公司已经拿到客户订单、合同价格或单品利润。


\begin{exhibitbox}[股东结构与机构持仓拥挤度（1）]
\scriptsize
\begin{tabularx}{\exhibitboxwidth}{L{0.8cm} L{1.05cm} L{1.3cm} L{0.75cm} L{0.75cm} L{0.8cm} L{0.9cm} L{0.75cm} X X}
\toprule
\textbf{代码} & \textbf{名称} & \textbf{股权结构} & \textbf{第一大} & \textbf{前十} & \textbf{流通前十} & \textbf{基金持仓} & \textbf{换手率} & \textbf{拥挤信号} & \textbf{模型动作} \\
\midrule
600549 & 厦门钨业 & 产业/国资控股约束强 & 28.4\% & 44.1\% & 49.7\% & 1只/0.0\% & 4.1\% & 中拥挤：高成交额；流通前十集中 & 维持市场锚，但提高公告和财报验证频率 \\
000657 & 中钨高新 & 产业/国资控股约束强 & 30.6\% & 66.1\% & 51.4\% & 1只/0.0\% & 4.8\% & 中拥挤：高成交额；流通前十集中 & 维持市场锚，但提高公告和财报验证频率 \\
002378 & 章源钨业 & 控股股东高度集中 & 56.7\% & 59.1\% & 61.9\% & 20只/0.3\% & 4.2\% & 低拥挤：流通前十集中 & 只作为流动性和机构参与度参考 \\
002842 & 翔鹭钨业 & 股权结构相对分散 & 15.8\% & 39.3\% & 27.2\% & 6只/0.2\% & 14.2\% & 中拥挤：高换手；流通前十中等集中 & 维持市场锚，但提高公告和财报验证频率 \\
603505 & 金石资源 & 控股股东高度集中 & 50.0\% & 63.4\% & 63.4\% & 201只/2.6\% & 2.9\% & 中拥挤：基金持仓中等；流通前十集中 & 维持市场锚，但提高公告和财报验证频率 \\
600160 & 巨化股份 & 控股股东高度集中 & 52.7\% & 57.4\% & 62.5\% & 139只/11.7\% & 3.7\% & 中拥挤：高成交额；基金持仓拥挤；流通前十集中 & 维持市场锚，但提高公告和财报验证频率 \\
603379 & 三美股份 & 前十股东集中 & 34.0\% & 65.9\% & 66.6\% & 82只/4.5\% & 2.6\% & 中拥挤：基金持仓中等；流通前十集中 & 维持市场锚，但提高公告和财报验证频率 \\
605020 & 永和股份 & 前十股东集中 & 36.5\% & 55.9\% & 56.4\% & 11只/3.2\% & 7.6\% & 中拥挤：中高换手；基金持仓中等；流通前十集中 & 维持市场锚，但提高公告和财报验证频率 \\
002407 & 多氟多 & 股权结构相对分散 & 10.3\% & 17.2\% & 13.0\% & 31只/2.8\% & 3.2\% & 低拥挤：基金持仓中等 & 只作为流动性和机构参与度参考 \\
\bottomrule
\end{tabularx}
\end{exhibitbox}

\begin{exhibitbox}[股东结构与机构持仓拥挤度（2）]
\scriptsize
\begin{tabularx}{\exhibitboxwidth}{L{0.8cm} L{1.05cm} L{1.3cm} L{0.75cm} L{0.75cm} L{0.8cm} L{0.9cm} L{0.75cm} X X}
\toprule
\textbf{代码} & \textbf{名称} & \textbf{股权结构} & \textbf{第一大} & \textbf{前十} & \textbf{流通前十} & \textbf{基金持仓} & \textbf{换手率} & \textbf{拥挤信号} & \textbf{模型动作} \\
\midrule
600378 & 昊华科技 & 产业/国资控股约束强 & 45.8\% & 73.4\% & 72.9\% & 16只/4.2\% & 5.0\% & 中拥挤：中高换手；基金持仓中等；流通前十集中 & 维持市场锚，但提高公告和财报验证频率 \\
300037 & 新宙邦 & 股权结构相对分散 & 13.7\% & 40.7\% & 21.0\% & 173只/10.9\% & 3.9\% & 低拥挤：基金持仓拥挤 & 只作为流动性和机构参与度参考 \\
688146 & 中船特气 & 控股股东高度集中 & 69.2\% & 80.4\% & 31.4\% & 7只/2.3\% & 10.6\% & 高拥挤：高换手；高成交额；低流通占比；基金持仓中等；流通前十中等集中 & 下调追价容忍度，市场锚需订单/业绩继续验证 \\
688549 & 中巨芯-U & 前十股东集中 & 26.4\% & 68.9\% & 29.2\% & 4只/0.9\% & 15.7\% & 高拥挤：高换手；低流通占比；流通前十中等集中 & 下调追价容忍度，市场锚需订单/业绩继续验证 \\
688268 & 华特气体 & 股权结构相对分散 & 0.0\% & 13.5\% & 56.0\% & 75只/2.4\% & 6.1\% & 中拥挤：中高换手；基金持仓中等；流通前十集中 & 维持市场锚，但提高公告和财报验证频率 \\
300346 & 南大光电 & 股权结构相对分散 & 9.8\% & 22.2\% & 23.4\% & 16只/2.2\% & 15.8\% & 中拥挤：高换手；高成交额；基金持仓中等 & 维持市场锚，但提高公告和财报验证频率 \\
688106 & 金宏气体 & 股权结构相对分散 & 20.2\% & 46.1\% & 49.9\% & 25只/1.8\% & 9.9\% & 中拥挤：中高换手；流通前十集中 & 维持市场锚，但提高公告和财报验证频率 \\
002549 & 凯美特气 & 股权结构相对分散 & 37.4\% & 45.5\% & 45.3\% & 37只/0.3\% & 8.5\% & 中拥挤：中高换手；流通前十集中 & 维持市场锚，但提高公告和财报验证频率 \\
002971 & 和远气体 & 前十股东集中 & 20.6\% & 54.8\% & 49.7\% & 54只/5.1\% & 13.8\% & 高拥挤：高换手；基金持仓拥挤；流通前十集中 & 下调追价容忍度，市场锚需订单/业绩继续验证 \\
\bottomrule
\end{tabularx}
\sourcenote{AkShare主要股东、流通股东和基金持股接口，叠加东方财富官方流动性快照；股东和基金持仓存在披露滞后，只用于二级市场拥挤度、机构参与度、市场锚权重和风险折扣，不替代客户订单、合同平均售价、收入确认比例、客户份额或单品毛利。完整快照见 data/shareholder\_crowding\_snapshot.json/md。}
\end{exhibitbox}

这组持仓数据把“换手很高”进一步拆成两类风险：若流通前十或基金持仓占比较高，短期成交放大更容易演变为拥挤交易；若股权集中度高但机构参与度低，价格弹性更多来自主题资金而非基本面长钱。估值模型因此只把它用于市场锚权重、追价容忍度和风险折扣；没有订单、公司合同平均售价、收入确认比例或单品毛利之前，基本面锚不能因为持仓拥挤而上修。


持有人结构之外，还要看杠杆资金是否参与推升。两融明细把“成交额放大”进一步拆成融资余额存量和融资买入流量：融资余额占流通市值越高，主题回撤时被动降杠杆压力越大；融资买入占成交额越高，当日价格越依赖杠杆资金持续流入。


\begin{exhibitbox}[融资融券杠杆拥挤度（1）]
\scriptsize
\begin{tabularx}{\exhibitboxwidth}{L{0.8cm} L{1.05cm} L{1.35cm} L{0.9cm} L{0.95cm} L{0.95cm} L{0.85cm} L{0.85cm} X X}
\toprule
\textbf{代码} & \textbf{名称} & \textbf{两融状态} & \textbf{数据日} & \textbf{融资余额} & \textbf{融资买入} & \textbf{融资/流通} & \textbf{买入/成交} & \textbf{杠杆信号} & \textbf{模型动作} \\
\midrule
600549 & 厦门钨业 & 两融标的 & 20260626 & 33.5亿元 & 7.8亿元 & 2.7\% & 14.5\% & 低杠杆压力：融资参与温和 & 只作为资金结构参考，不调整基本面目标价 \\
000657 & 中钨高新 & 两融标的 & 20260626 & 33.3亿元 & 8.8亿元 & 2.5\% & 13.0\% & 低杠杆压力：融资参与温和 & 只作为资金结构参考，不调整基本面目标价 \\
002378 & 章源钨业 & 两融标的 & 20260626 & 14.7亿元 & 2.8亿元 & 3.6\% & 16.4\% & 低杠杆压力：融资余额占比中等；融资参与温和 & 只作为资金结构参考，不调整基本面目标价 \\
002842 & 翔鹭钨业 & 非两融/未覆盖 & 20260626 & 0.0亿元 & 0.0亿元 & 0.0\% & 0.0\% & 非两融标的：当日明细未覆盖 & 不作为杠杆拥挤交易，但保留流动性观察 \\
603505 & 金石资源 & 两融标的 & 20260626 & 6.4亿元 & 0.8亿元 & 4.1\% & 16.4\% & 低杠杆压力：融资余额占比中等；融资参与温和 & 只作为资金结构参考，不调整基本面目标价 \\
600160 & 巨化股份 & 两融标的 & 20260626 & 46.7亿元 & 7.4亿元 & 3.4\% & 14.6\% & 低杠杆压力：融资余额占比中等；融资参与温和 & 只作为资金结构参考，不调整基本面目标价 \\
603379 & 三美股份 & 两融标的 & 20260626 & 12.9亿元 & 1.0亿元 & 3.2\% & 9.2\% & 低杠杆压力：融资余额占比中等 & 只作为资金结构参考，不调整基本面目标价 \\
605020 & 永和股份 & 两融标的 & 20260626 & 6.5亿元 & 1.7亿元 & 3.2\% & 10.8\% & 低杠杆压力：融资余额占比中等；融资参与温和 & 只作为资金结构参考，不调整基本面目标价 \\
002407 & 多氟多 & 两融标的 & 20260626 & 29.1亿元 & 7.1亿元 & 5.9\% & 44.3\% & 高杠杆拥挤：融资余额占比较高；融资买入占成交额高 & 下调追价容忍度，优先等公告/财报验证 \\
\bottomrule
\end{tabularx}
\end{exhibitbox}

\begin{exhibitbox}[融资融券杠杆拥挤度（2）]
\scriptsize
\begin{tabularx}{\exhibitboxwidth}{L{0.8cm} L{1.05cm} L{1.35cm} L{0.9cm} L{0.95cm} L{0.95cm} L{0.85cm} L{0.85cm} X X}
\toprule
\textbf{代码} & \textbf{名称} & \textbf{两融状态} & \textbf{数据日} & \textbf{融资余额} & \textbf{融资买入} & \textbf{融资/流通} & \textbf{买入/成交} & \textbf{杠杆信号} & \textbf{模型动作} \\
\midrule
600378 & 昊华科技 & 两融标的 & 20260626 & 19.7亿元 & 6.6亿元 & 3.1\% & 19.3\% & 低杠杆压力：融资余额占比中等；融资参与温和 & 只作为资金结构参考，不调整基本面目标价 \\
300037 & 新宙邦 & 两融标的 & 20260626 & 15.7亿元 & 3.7亿元 & 3.5\% & 21.0\% & 中杠杆拥挤：融资余额占比中等；融资买入活跃 & 保留市场锚，但提高杠杆回撤监测频率 \\
688146 & 中船特气 & 两融标的 & 20260626 & 30.6亿元 & 10.3亿元 & 6.2\% & 19.1\% & 中杠杆拥挤：融资余额占比较高；融资参与温和；当日融资净买入 & 保留市场锚，但提高杠杆回撤监测频率 \\
688549 & 中巨芯-U & 两融标的 & 20260626 & 10.8亿元 & 3.0亿元 & 5.3\% & 9.4\% & 低杠杆压力：融资余额占比较高 & 只作为资金结构参考，不调整基本面目标价 \\
688268 & 华特气体 & 两融标的 & 20260626 & 14.5亿元 & 4.1亿元 & 4.7\% & 21.3\% & 中杠杆拥挤：融资余额占比中等；融资买入活跃 & 保留市场锚，但提高杠杆回撤监测频率 \\
300346 & 南大光电 & 两融标的 & 20260626 & 31.8亿元 & 9.0亿元 & 5.6\% & 9.9\% & 低杠杆压力：融资余额占比较高 & 只作为资金结构参考，不调整基本面目标价 \\
688106 & 金宏气体 & 两融标的 & 20260626 & 10.8亿元 & 2.5亿元 & 5.3\% & 12.5\% & 中杠杆拥挤：融资余额占比较高；融资参与温和 & 保留市场锚，但提高杠杆回撤监测频率 \\
002549 & 凯美特气 & 两融标的 & 20260626 & 4.2亿元 & 1.2亿元 & 2.9\% & 9.8\% & 低杠杆压力：融资信号温和 & 只作为资金结构参考，不调整基本面目标价 \\
002971 & 和远气体 & 非两融/未覆盖 & 20260626 & 0.0亿元 & 0.0亿元 & 0.0\% & 0.0\% & 非两融标的：当日明细未覆盖 & 不作为杠杆拥挤交易，但保留流动性观察 \\
\bottomrule
\end{tabularx}
\sourcenote{上交所融资融券明细、深交所融资融券Excel公开接口，叠加东方财富官方流通市值和成交额；两融明细只用于杠杆资金参与度、回撤放大风险和市场锚压力，不替代客户订单、合同平均售价、收入确认比例、客户份额或单品毛利。完整快照见 data/margin\_leverage\_snapshot.json/md。}
\end{exhibitbox}

这组两融数据的作用是约束“能不能追价”，不是证明“值不值这个价”。如果融资买入占成交额已经很高，而公司仍没有订单金额、合同平均售价、收入确认比例和单品毛利，模型应降低追价容忍度；若融资信号温和，则说明短期回撤主要来自普通流动性和主题情绪，不构成基本面上修证据。


龙虎榜把二级市场证据再推进一层：它只披露满足异常交易条件后的公开席位和机构净买卖，可以判断主题资金是否已经进入高频上榜、机构净买入或机构净卖出状态。未上榜不代表没有资金参与，上榜也不代表订单、价格或毛利已经兑现；因此它只能影响市场锚权重、追价容忍度和回撤风险，不能提高基本面目标价。


\begin{exhibitbox}[龙虎榜席位异常与机构交易（1）]
\scriptsize
\begin{tabularx}{\exhibitboxwidth}{L{0.8cm} L{1.05cm} L{1.2cm} L{0.9cm} L{0.55cm} L{0.9cm} L{0.9cm} L{0.9cm} X X}
\toprule
\textbf{代码} & \textbf{名称} & \textbf{近一月状态} & \textbf{最近上榜} & \textbf{次数} & \textbf{净买额} & \textbf{总成交} & \textbf{机构净买} & \textbf{席位信号} & \textbf{模型动作} \\
\midrule
600549 & 厦门钨业 & 近一月上榜 & 2026-06-18 & 3 & 18.3亿元 & 177.3亿元 & 20.7亿元 & 高席位拥挤：多次上榜；龙虎榜净额波动大；机构净买入 & 下调追价容忍度，优先等待公告/财报验证 \\
000657 & 中钨高新 & 近一月上榜 & 2026-06-22 & 2 & 10.1亿元 & 52.2亿元 & 1.6亿元 & 中席位拥挤：多次上榜；龙虎榜净额波动大；机构分歧 & 保留市场锚，但提高席位和回撤监测频率 \\
002378 & 章源钨业 & 近一月上榜 & 2026-06-17 & 2 & 2.3亿元 & 17.3亿元 & 0.4亿元 & 低席位扰动：多次上榜 & 只作为交易扰动参考，不调整基本面目标价 \\
002842 & 翔鹭钨业 & 近一月上榜 & 2026-06-26 & 8 & 3.3亿元 & 91.8亿元 & 3.8亿元 & 高席位拥挤：高频上榜；龙虎榜净额显著；机构净买入；6月明细密集 & 下调追价容忍度，优先等待公告/财报验证 \\
603505 & 金石资源 & 近一月未上榜 & 未上榜 & 0 & 0.0亿元 & 0.0亿元 & 0.0亿元 & 近一月未上榜：无公开异常席位信号 & 无公开异常席位，不作为交易热度加分 \\
600160 & 巨化股份 & 近一月未上榜 & 未上榜 & 0 & 0.0亿元 & 0.0亿元 & 0.0亿元 & 近一月未上榜：无公开异常席位信号 & 无公开异常席位，不作为交易热度加分 \\
603379 & 三美股份 & 近一月未上榜 & 未上榜 & 0 & 0.0亿元 & 0.0亿元 & 0.0亿元 & 近一月未上榜：无公开异常席位信号 & 无公开异常席位，不作为交易热度加分 \\
605020 & 永和股份 & 近一月未上榜 & 未上榜 & 0 & 0.0亿元 & 0.0亿元 & 0.0亿元 & 近一月未上榜：无公开异常席位信号 & 无公开异常席位，不作为交易热度加分 \\
002407 & 多氟多 & 近一月上榜 & 2026-06-22 & 2 & 6.6亿元 & 30.4亿元 & -3.3亿元 & 中席位拥挤：多次上榜；龙虎榜净额显著；机构净卖出 & 保留市场锚，但提高席位和回撤监测频率 \\
\bottomrule
\end{tabularx}
\end{exhibitbox}

\begin{exhibitbox}[龙虎榜席位异常与机构交易（2）]
\scriptsize
\begin{tabularx}{\exhibitboxwidth}{L{0.8cm} L{1.05cm} L{1.2cm} L{0.9cm} L{0.55cm} L{0.9cm} L{0.9cm} L{0.9cm} X X}
\toprule
\textbf{代码} & \textbf{名称} & \textbf{近一月状态} & \textbf{最近上榜} & \textbf{次数} & \textbf{净买额} & \textbf{总成交} & \textbf{机构净买} & \textbf{席位信号} & \textbf{模型动作} \\
\midrule
600378 & 昊华科技 & 近一月上榜 & 2026-06-17 & 3 & -10.7亿元 & 69.2亿元 & -5.8亿元 & 高席位拥挤：多次上榜；龙虎榜净额波动大；机构净卖出 & 下调追价容忍度，优先等待公告/财报验证 \\
300037 & 新宙邦 & 近一月未上榜 & 未上榜 & 0 & 0.0亿元 & 0.0亿元 & 0.0亿元 & 近一月未上榜：无公开异常席位信号 & 无公开异常席位，不作为交易热度加分 \\
688146 & 中船特气 & 近一月上榜 & 2026-06-11 & 4 & -11.0亿元 & 577.3亿元 & 0.7亿元 & 中席位拥挤：多次上榜；龙虎榜净额波动大 & 保留市场锚，但提高席位和回撤监测频率 \\
688549 & 中巨芯-U & 近一月上榜 & 2026-06-12 & 8 & -12.0亿元 & 292.7亿元 & -3.2亿元 & 高席位拥挤：高频上榜；龙虎榜净额波动大；机构净卖出；6月明细密集 & 下调追价容忍度，优先等待公告/财报验证 \\
688268 & 华特气体 & 近一月上榜 & 2026-06-23 & 5 & 1.7亿元 & 48.2亿元 & -0.3亿元 & 中席位拥挤：高频上榜；6月明细密集 & 保留市场锚，但提高席位和回撤监测频率 \\
300346 & 南大光电 & 近一月未上榜 & 未上榜 & 0 & 0.0亿元 & 0.0亿元 & 0.0亿元 & 近一月未上榜：无公开异常席位信号 & 无公开异常席位，不作为交易热度加分 \\
688106 & 金宏气体 & 近一月未上榜 & 未上榜 & 0 & 0.0亿元 & 0.0亿元 & 0.0亿元 & 近一月未上榜：无公开异常席位信号 & 无公开异常席位，不作为交易热度加分 \\
002549 & 凯美特气 & 近一月未上榜 & 未上榜 & 0 & 0.0亿元 & 0.0亿元 & 0.0亿元 & 近一月未上榜：无公开异常席位信号 & 无公开异常席位，不作为交易热度加分 \\
002971 & 和远气体 & 近一月上榜 & 2026-06-16 & 6 & -3.8亿元 & 48.7亿元 & -2.0亿元 & 高席位拥挤：高频上榜；龙虎榜净额显著；机构分歧；6月明细密集 & 下调追价容忍度，优先等待公告/财报验证 \\
\bottomrule
\end{tabularx}
\sourcenote{东方财富龙虎榜公开接口，含近一月个股上榜统计、机构席位统计和6月明细；龙虎榜只用于公开异常交易席位、机构净买卖、追价容忍度和回撤风险监测，不替代客户订单、合同平均售价、收入确认比例、客户份额或单品毛利。完整快照见 data/dragon\_tiger\_seat\_snapshot.json/md。}
\end{exhibitbox}

这组席位数据的投资含义是“不要把短期机构交易误读为订单兑现”。高席位拥挤标的应降低追价容忍度并等待公告或财报验证；中席位拥挤标的可以保留市场锚，但必须提高席位和回撤监测频率；未上榜标的则不能因为没有龙虎榜记录就排除资金参与，只能说明公开异常席位证据不足。


\begin{exhibitbox}[交易风险公告与估值热度]
\scriptsize
\begin{tabularx}{\exhibitboxwidth}{L{0.9cm} L{1.2cm} X X X}
\toprule
\textbf{代码} & \textbf{名称} & \textbf{交易热度/估值} & \textbf{官方边界} & \textbf{估值用途} \\
\midrule
688146 & 中船特气 & 2026-05-11至2026-06-25累计涨幅388.77\%；滚动市盈率 573.30x；静态市盈率 597.55x；行业近一个月平均滚动市盈率 65.70x；静态市盈率 71.85x & WF6现有产能2000吨/年、6N级；未公开披露产品价格；订单达到披露标准时依法披露；上游钨材料上涨可能挤压利润空间。 & 高拥挤/高估值风险锚；市场锚权重必须由后续订单、价格和毛利验证维持。 \\
688549 & 中巨芯-U & 2026-05-08至2026-06-18累计涨幅180.63\%；2025年归母净利润为负，传统市盈率锚弱；C39行业滚动市盈率 71.45x；静态市盈率 78.31x & 高纯WF6未签署新的具有法律约束力的长期或大额实质性订单；产能600吨；暂无扩产计划；钨材料涨价带来成本不确定性。 & 高涨幅与订单缺口并存，市销率/现金流观察权重高于每股收益信用。 \\
600378 & 昊华科技 & 2026-06-04至2026-06-11连续6个交易日累计涨幅53.40\%；公告未给出市盈率；报告使用本报告估值表约束；公告未给出行业市盈率；报告使用同业估值与主营构成约束 & 2025年度WF6收入占营业收入0.13\%，占比较低，不会对公司业绩产生重大影响；不存在应披露未披露重大信息。 & WF6题材热度不得覆盖主营平台估值，市场锚需向平台业务和分部估值回归。 \\
002971 & 和远气体 & 2026-06-04、06-05、06-08连续3个交易日累计偏离21.42\%；公告未给出市盈率；报告使用本报告估值表约束；公告未给出行业市盈率；报告使用同业估值与项目阶段约束 & 电子级WF6等产品尚处试生产阶段；WF6尚未产生任何业绩，未签署任何具有法律约束力的实质性订单协议。 & 项目期权不能转为当期增长每股收益，市场锚应低权重处理。 \\
\bottomrule
\end{tabularx}
\sourcenote{公司异常波动、风险提示和停牌核查公告PDF归档整理；交易风险公告只约束市场锚和风险折扣，不替代订单、合同平均售价、收入确认比例或单品毛利。}
\end{exhibitbox}

这张表把市场热度转成了可执行约束：中船特气和中巨芯-U已经出现显著超越指数的累计涨幅，且中船特气公告口径下滚动市盈率达到573.30倍；昊华科技和和远气体则分别明确WF6收入占比低、试生产尚未产生业绩。对投资结论的含义不是简单看空，而是把市场锚权重降到“需要下一份订单、价格、毛利或收入确认继续验证”的层级。若后续公告仍停留在产能、洽谈或试生产，报告的目标价只能维持基本面锚，不能因为交易热度自动上修。
